\documentclass[10pt,a4paper]{amsart}
\usepackage[margin=19mm]{geometry}
\usepackage{amsmath,amssymb,mathtools}
\usepackage[scr=rsfs,cal=euler]{mathalfa}
\usepackage{xfrac}
\usepackage[unicode,hidelinks,bookmarks=false]{hyperref}
\newcommand{\ie}{i.e.\ }
\newcommand{\cf}{cf.\ }
\newcommand{\resp}{respectively\ }
\newcommand{\Subsection}{\subsection}
\providecommand{\nobreakdash}{\nobreak\hbox{-}}
\setlength{\emergencystretch}{2em}
\multlinegap0pt
\usepackage{bm}
\usepackage{bbm}
\usepackage{dsfont}
\usepackage{xspace}
\usepackage{cancel}
\usepackage{mathtools}
\usepackage{amsthm}
\makeatletter
\@ifundefined{satz}{\newtheorem{satz}{Satz}}{}
\@ifundefined{assumptions}{\newtheorem{assumptions}[satz]{Assumptions}}{}
\makeatother
\DeclareMathOperator	{\IE}		{\mathbb{E}}
\DeclareMathOperator	{\IN}		{\mathbb{N}}
\DeclarePairedDelimiter	{\norm}		{\lVert}	{\rVert}
\def\E{{\mathbb E}}
\def\R{{\mathbb R}}
\title[Some notes on higher order concentration of measure]{Some notes on higher order concentration of measure}
\author{Holger Sambale}
\date{Source: arXiv:2507.08692v1 (CC BY 4.0)}
\begin{document}
\maketitle

\noindent\emph{Source and licence.} Some notes on higher order concentration of measure. Version arXiv:2507.08692v1 (CC BY 4.0), distributed under the Creative Commons Attribution 4.0 International licence, as recorded in the arXiv licence metadata for this exact version and shown on the abstract page https://arxiv.org/abs/2507.08692v1. Full rights notice: licenses/arxiv-2507.08692-CC-BY-4.0.txt.\par\smallskip

\noindent\textit{Editorial scope.} Counted unit: the Satz whose source label is \texttt{thm:HigherOrder2} in the article (source0.tex lines 1040--1061, source label \texttt{thm:HigherOrder2}), delivered together with the article's own complete original proof of it (source0.tex lines 1075--1097, one \texttt{proof} environment of 23 lines). No mathematics is written, repaired or re-derived: statement and proof are the article's own text line for line. The proof is detached from the statement in the source: the article states the result at lines 1040--1061 and prints its proof at lines 1075--1097, and the proof environment's own header names the statement, so the association is the article's own. Required context retained uncounted: ass:Set1 (lines 208--221); thm:HigherOrder (lines 229--250); ass:Set2 (lines 1014--1036). Every definition, display and label of those lines is retained unchanged. Omitted from this package and not redistributed: 1--207, 222--228, 251--1013, 1037--1039, 1062--1074, 1098--1611. Nothing inside a retained excerpt is omitted or altered: every retained body is delivered exactly as the article's lines are written, so no omission receipt of any kind accompanies this package. Citations: the retained excerpts carry no citation command at all, so no bibliography entry of the article is transcribed and no reference is redistributed. Reference adaptation: none was needed and none was made; every reference inside the delivered unit resolves inside it, so no unresolved reference or citation is produced. Printed slips kept literal: no printed word, symbol, hypothesis or number of the counted statement or of its proof was altered. Layout and class: the article's own class is \texttt{amsart} with packages including inputenc, babel, fontenc, csquotes, amssymb, amsthm, mathtools, bbm, ifxetex, mathrsfs, trfsigns, xspace, geometry, enumerate; the wrapper is the lane's pinned \texttt{amsart} editorial wrapper at 10pt on A4, which restates the theorem-like environments the retained text needs and adds the article's own macro definitions that the retained text needs (\texttt{\textbackslash E}, \texttt{\textbackslash IE}, \texttt{\textbackslash IN}, \texttt{\textbackslash R}, \texttt{\textbackslash norm}), with extra wrapper packages (bm, bbm, dsfont, xspace, cancel, mathtools). The delivered numbering is the unit's own. Assets: no retained span contains a figure environment, a graphics-inclusion command, a PDF-page inclusion, a table or a TikZ picture; no image file is copied into this package, the package declares no asset dependency and no image file is redistributed with it. Licence: version arXiv:2507.08692v1 of this article is distributed under the Creative Commons Attribution 4.0 International licence, as recorded in the arXiv licence metadata for this exact version and shown on the abstract page https://arxiv.org/abs/2507.08692v1. A full rights notice is delivered at licenses/arxiv-2507.08692-CC-BY-4.0.txt. Characters: every retained excerpt is pure ASCII, so the delivered engine, XeLaTeX with its default Latin Modern text font, reports no missing character.
\begin{assumptions}\label{ass:Set1}
\begin{enumerate}
\item There exist $p > 0$, $r_0 > 1$, and $L, \sigma > 0$ such that for any function $g \in \mathcal{F}$, we have
\begin{equation*}
	\lVert g - \E g \rVert_r \le L\sigma r^{1/p} \lVert \Gamma g \rVert_r
\end{equation*}
for all $r \ge r_0$.
\item There exist ``higher order'' difference operators $\Gamma^{(j)}(g) \in [0,\infty)$, $j = 1, \ldots, d$, such that for any $g \in \mathcal{F}$,
\begin{equation*}
	\Gamma (\Gamma^{(j)}(g)) \le \Gamma^{(j+1)}(g)
\end{equation*}
for any $j = 1, \ldots, d-1$.
\end{enumerate}
\end{assumptions}

\begin{satz}\label{thm:HigherOrder}
    Given Assumptions \ref{ass:Set1}, let $f\colon \mathcal{X} \to \R$ be a suitable function with $\E f = 0$.
    \begin{enumerate}
        \item Assuming that $\lVert \Gamma^{(j)} f \rVert_1 \le \sigma^{d-j}$ for all $j=1,\ldots,d-1$ and $\lVert \Gamma^{(d)} f \rVert_\infty \le 1$, it holds that
        \[
        \int \exp\Big\{\frac{c_{p,d,r_0,L}}{\sigma^p} |f|^{p/d} \Big\} d\mu \le 2,
        \]
        where a possible choice of the constant $c_{p,d,r_0,L}$ is given by
        \[
        c_{p,d,r_0,L} = \frac{(r_0^{1/p}-1)^p}{2e\max(L^{1/d}, L)^pr_0\max(r_0,p/d)}.
        \]
        \item For any $t\ge 0$, it holds that
        \[
        \mu(|f| \ge t)\le
        2 \exp\Big\{- \frac{C_{p,r_0,L}}{d^p\sigma^p} \min\Big(\min_{j=1, \ldots,d-1} \Big(\frac{t}{\lVert \Gamma^{(j)} f \rVert_1}\Big)^{p/j}, \Big(\frac{t}{\lVert \Gamma^{(d)} f \rVert_\infty}\Big)^{p/d}\Big)\Big\},
        \]
    where a possible choice of the constant $C_{p,r_0,L}$ is given by
    \[
    C_{p,r_0,L} = \frac{\log(2)}{r_0(Le)^p}.
    \]
    \end{enumerate}
\end{satz}

\begin{assumptions}\label{ass:Set2}
	\begin{enumerate}
		\item There exist $p > 0$, $r_0 > 1$, and $L, \sigma > 0$ such that for any function $g \in \mathcal{F}$, we have
		\begin{equation*}%\label{eq:Cond1}
			\lVert g - \E g \rVert_r \le L\sigma r^{1/p} \lVert \Gamma g \rVert_r
		\end{equation*}
		for all $r \ge r_0$.
		\item For the same quantities $p > 0$, $r_0 > 1$ and $L, \sigma > 0$ as in (1) and any function $g \in \mathcal{F}$, we have
		\begin{equation*}%\label{eq:Cond1}
			\lVert (g - \E g)_+ \rVert_r \le L\sigma r^{1/p} \lVert \Gamma^+ g \rVert_r
		\end{equation*}
		\item There exist ``higher order'' difference operators $\Gamma^{(j)}(g) \in [0,\infty)$, $j = 1, \ldots, d$, such that for any $g \in \mathcal{F}$,
		\begin{equation*}%\label{eq:Cond2}
			\Gamma^+ (\Gamma^{(j)}(g)) \le \Gamma^{(j+1)}(g)
		\end{equation*}
		for any $j = 1, \ldots, d-1$.
		\item[(3')] There exist ``higher order'' difference operators $\Gamma^{(j)}(g) \in [0,\infty)$, $j = 1, \ldots, d$, and some constant $\gamma > 0$ such that for any $g \in \mathcal{F}$, $\Gamma^+(g) \le \gamma \Gamma(g)$ as well as
		\begin{equation*}%\label{eq:Cond2}
			\Gamma^+ (\Gamma^{(j)}(g)) \le \gamma \Gamma^{(j+1)}(g)
		\end{equation*}
		for any $j = 1, \ldots, d-1$.
	\end{enumerate}
\end{assumptions}

\begin{satz}\label{thm:HigherOrder2}
	Given Assumptions \ref{ass:Set2} (1)--(3), let $f\colon \Omega \to \R$ be a suitable function with $\E f = 0$.
	\begin{enumerate}
		\item Assuming that $\lVert \Gamma^{(j)} f \rVert_1 \le \sigma^{d-j}$ for all $j=1,\ldots,d-1$ and $\lVert \Gamma^{(d)} f \rVert_\infty \le 1$, it holds that
		\[
		\int \exp\Big\{\frac{c_{p,d,r_0,L}}{\sigma^p} |f|^{p/d} \Big\} d\mu \le 2,
		\]
		where a possible choice of the constant $c_{p,d,r_0,L}$ is given by
		\[
		c_{p,d,r_0,L} = \frac{(r_0^{1/p}-1)^p}{2e\max(L^{1/d}, L)^pr_0\max(r_0,p/d)}.
		\]
		\item For any $t\ge 0$, it holds that
		\[
		\mu(|f| \ge t)\le
		2 \exp\Big\{- \frac{C_{p,r_0,L}}{(d\sigma)^p} \min\Big(\min_{k=1, \ldots,d-1} \Big(\frac{t}{\lVert \Gamma^{(k)} f \rVert_1}\Big)^{p/k}, \Big(\frac{t}{\lVert \Gamma^{(d)} f \rVert_\infty}\Big)^{p/d}\Big)\Big\},
		\]
		where a possible choice of the constant $C_{p,r_0,L}$ is given by
		\[
	    C_{p,r_0,L} = \frac{\log(2)}{r_0(Le)^p}.
		\]
	\end{enumerate}
\end{satz}

\begin{proof}[Proof of Theorem \ref{thm:HigherOrder2}]
	The proof is very similar to the proof of Theorem \ref{thm:HigherOrder}. First, Assumption \ref{ass:Set2} (1) yields
	\begin{align*}
		\norm{f}_r \le L\sigma r^{1/p} \norm{\Gamma f}_p \le L\sigma r^{1/p} \norm{\Gamma f}_1 + L\sigma r^{1/p} \norm{(\Gamma f - \E \Gamma f)_+}_r
	\end{align*}
	where we have used that for any positive random variable $W$
	\begin{align}		\label{eqn:LpEstimatePositiveFunction}
		\norm{W}_r \le \IE W + \norm{(W - \IE W)_+}_r.
	\end{align}
	Now, by Assumptions \ref{ass:Set2} (2)\&(3) it follows that
	\begin{align*}
		\norm{\left( \Gamma f - \E \Gamma f \right)_+}_r \le L\sigma r^{1/p} \norm{\Gamma^+ (\Gamma f)}_r \le L\sigma r^{1/p} \norm{\Gamma^{(2)}f}_r.
	\end{align*}
	This can be easily iterated to obtain for any $d \in \IN$
	\begin{align*}
		\norm{f}_r \le \sum_{j = 1}^{d-1} (L\sigma r^{1/p})^j \norm{\Gamma^{(j)}f}_1 + (L\sigma r^{1/p})^d \norm{\Gamma^{(d)}f}_\infty.
	\end{align*}
	If (1) is not satisfied, then one uses (2) in the first step instead, which leads to
	\begin{align*}
		\norm{f_+}_r \le \sum_{j = 1}^{d-1} (L\sigma r^{1/p})^j \norm{\Gamma^{(j)}f}_1 + (L\sigma r^{1/p})^d \norm{\Gamma^{(d)}f}_\infty.
	\end{align*}
	Now one can argue as in the proof of Theorem \ref{thm:HigherOrder}. The case of (3) being replaced by (3') follows in the same way.
\end{proof}

\end{document}
